\documentclass{article}

\usepackage{amsthm}% para usar o comando \theoremstyle que especifica a formatação escolhida para o ambiente theorem

\newtheorem{teo}{Teorema}
	
\theoremstyle{definition}%define a formatação do teorema
\newtheorem{ex}[teo]{Exemplo}

\begin{document}

\begin{teo}
    asndkasdaskdjkj
\end{teo}

\begin{ex}
    asdasdasdasdas
\end{ex}


$\mathcal{A}$
$\mathcal{V}$

\vspace{.5cm}

\begin{minipage}{.5\textwidth}
    aaaaaaaaaaaaaaaaaaaaaaaaaaaa aaaaaaaaaaaaaaaaaaaaaaaaaaaa aaaaaaaaaaaaaaaaaaaaaaaaaaaa aaaaaaaaaaaaaaaaaaaaaaaaaaaa aaaaaaaaaaaaaaaaaaaaaaaaaaaa aaaaaaaaaaaaaaaaaaaaaaaaaaaa aaaaaaaaaaaaaaaaaaaaaaaaaaaa aaaaaaaaaaaaaaaaaaaaaaaaaaaa aaaaaaaaaaaaaaaaaaaaaaaaaaaa aaaaaaaaaaaaaaaaaaaaaaaaaaaa aaaaaaaaaaaaaaaaaaaaaaaaaaaa aaaaaaaaaaaaaaaaaaaaaaaaaaaa aaaaaaaaaaaaaaaaaaaaaaaaaaaa aaaaaaaaaaaaaaaaaaaaaaaaaaaa aaaaaaaaaaaaaaaaaaaaaaaaaaaa aaaaaaaaaaaaaaaaaaaaaaaaaaaa aaaaaaaaaaaaaaaaaaaaaaaaaaaa aaaaaaaaaaaaaaaaaaaaaaaaaaaa aaaaaaaaaaaaaaaaaaaaaaaaaaaa aaaaaaaaaaaaaaaaaaaaaaaaaaaa aaaaaaaaaaaaaaaaaaaaaaaaaaaa aaaaaaaaaaaaaaaaaaaaaaaaaaaa aaaaaaaaaaaaaaaaaaaaaaaaaaaa aaaaaaaaaaaaaaaaaaaaaaaaaaaa
\end{minipage}
\begin{minipage}{.5\textwidth}
    aaaaaaaaaaaaaaaaaaaaaaaaaaaa aaaaaaaaaaaaaaaaaaaaaaaaaaaa aaaaaaaaaaaaaaaaaaaaaaaaaaaa aaaaaaaaaaaaaaaaaaaaaaaaaaaa aaaaaaaaaaaaaaaaaaaaaaaaaaaa aaaaaaaaaaaaaaaaaaaaaaaaaaaa aaaaaaaaaaaaaaaaaaaaaaaaaaaa aaaaaaaaaaaaaaaaaaaaaaaaaaaa aaaaaaaaaaaaaaaaaaaaaaaaaaaa aaaaaaaaaaaaaaaaaaaaaaaaaaaa aaaaaaaaaaaaaaaaaaaaaaaaaaaa aaaaaaaaaaaaaaaaaaaaaaaaaaaa aaaaaaaaaaaaaaaaaaaaaaaaaaaa aaaaaaaaaaaaaaaaaaaaaaaaaaaa aaaaaaaaaaaaaaaaaaaaaaaaaaaa aaaaaaaaaaaaaaaaaaaaaaaaaaaa aaaaaaaaaaaaaaaaaaaaaaaaaaaa aaaaaaaaaaaaaaaaaaaaaaaaaaaa aaaaaaaaaaaaaaaaaaaaaaaaaaaa aaaaaaaaaaaaaaaaaaaaaaaaaaaa aaaaaaaaaaaaaaaaaaaaaaaaaaaa aaaaaaaaaaaaaaaaaaaaaaaaaaaa aaaaaaaaaaaaaaaaaaaaaaaaaaaa aaaaaaaaaaaaaaaaaaaaaaaaaaaa
\end{minipage}


\end{document}